\subsubsection{Space-Time Numbers}
A spacetime increment is the real paravector
\[
\dd\pv{X}=\dd t+\dd r_{1}\pv{\sigma}_{1}+\dd r_{2}\pv{\sigma}_{2}+\dd r_{3}\pv{\sigma}_{3}
=\dd t+\dd\vct{r}\cdot\sigv,
\]
with conjugate
\[
\conj{\dd\pv{X}}=\dd t-\dd r_{1}\pv{\sigma}_{1}-\dd r_{2}\pv{\sigma}_{2}-\dd r_{3}\pv{\sigma}_{3}
=\dd t-\dd\vct{r}\cdot\sigv.
\]
Here \(\dd t\in\R\) is the scalar part and \(\dd\vct{r}\in\R^{3}\) is the spatial vector part. Their recovery formulas are
\begin{equation*}
\frac{\dd\pv{X}+\conj{\dd\pv{X}}}{2}=\dd t,
\qquad
\frac{\dd\pv{X}-\conj{\dd\pv{X}}}{2}=\dd\vct{r}\cdot\sigv,
\end{equation*}
so
\begin{equation*}
\eqtext{scalar}(\dd\pv{X})=\dd t,
\qquad
\eqtext{vector}(\dd\pv{X})=\dd\vct{r}.
\end{equation*}
The determinant of the increment is
\begin{equation*}
\detf{\dd\pv{X}}
=\dd\pv{X}\conj{\dd\pv{X}}
=\conj{\dd\pv{X}}\dd\pv{X}
=\dd t^{2}-\dd\vct{r}^{2}
=\dd s^{2}.
\end{equation*}
The familiar Lorentzian identities are obeyed:
\begin{align*}
\dd t^2>\dd\vct{r}^2,\ \dd\pv{X}\neq 0
&\eqtext{ implies }\detf{\dd\pv{X}}>0,\notag\\
\detf{\dd\pv{X}+\dd\pv{X}'}
&-\detf{\dd\pv{X}}\notag\\
&\qquad -\detf{\dd\pv{X}'}
=2(\dd t\,\dd t'-\dd\vct{r}\cdot\dd\vct{r}'),\notag\\
\detf{a\,\dd\pv{X}}
&=a^2\detf{\dd\pv{X}},
\qquad a\in\R.
\end{align*}
The commutator of two spacetime increments is purely spatial:
\begin{align*}
\com{\dd\pv{X}}{\dd\pv{X}'}
&=\com{\dd\vct{r}\cdot\sigv}{\dd\vct{r}'\cdot\sigv}\notag\\
&=2i(\dd\vct{r}\times\dd\vct{r}')\cdot\sigv.
\end{align*}
Along a future-directed time-like path, so that \(\dd t>0\), the following relation is satisfied by the positive arc length:
\begin{align*}
\dd s
&=\sqrt{\dd t^{2}-\dd\vct{r}^{2}}\notag\\
&=\sqrt{1-\dd\vct{r}^{2}/\dd t^{2}}\,\dd t\notag\\
&=\sqrt{1-\vct{v}^{2}}\,\dd t,
\end{align*}
where
\begin{equation*}
\vct{v}:=\frac{\dd\vct{r}}{\dd t}=(v_{1},v_{2},v_{3})\in\R^{3},
\qquad
\gamma:=\frac{\dd t}{\dd s}=\frac{1}{\sqrt{1-\vct{v}^{2}}}.
\end{equation*}
For a real non-null increment, \(\herm{\dd\pv{X}}=\dd\pv{X}\), so the inverse exists and
\begin{align*}
(\dd\pv{X})^{-1}
&=\frac{\adjf{\dd\pv{X}}}{\detf{\dd\pv{X}}}\notag\\
&=\frac{\conj{\dd\pv{X}}}{\detf{\dd\pv{X}}}\notag\\
&=\frac{\conj{\dd\pv{X}}}{\dd\pv{X}\conj{\dd\pv{X}}}\notag\\
&=\frac{\conj{\dd\pv{X}}}{\dd s^{2}}.
\end{align*}
Null displacements are singled out algebraically because this inverse breaks down exactly when \(\dd s^{2}=0\).

\subsubsection{Interval Classes and Proper Time}
For a spacetime increment \(\dd\pv{X}=\dd t+\dd \vct{r}\cdot\sigv\), the interval is
\begin{equation*}
\detf{\dd\pv{X}}=\dd\pv{X}\,\conj{\dd\pv{X}}=\dd t^2-\dd \vct{r}^2=\dd s^2.
\end{equation*}
The three standard interval classes are obtained from this identity without any additional geometric machinery:
\begin{align*}
\dd s^2>0 &\quad\eqtext{implies a time-like interval},\notag \\
\dd s^2=0,\ \dd\pv{X}\neq0
&\quad\eqtext{implies a null interval},\notag \\
\dd s^2<0 &\quad\eqtext{implies a space-like interval}.
\end{align*}
If \(\detf{\dd\pv{X}}>0\), the interval is time-like. Then the arc length \(s\) equals the proper time \(\tau\) measured along the path:
\begin{align*}
\tau=s=\int \dd s=\int \dd\tau
&=\int \sqrt{1-\vct{v}^2}\,\dd t\notag\\
&=\int \frac{1}{\gamma}\,\dd t.
\end{align*}
This is the time-dilation relation in natural units. In SI units the same statement is
\begin{equation*}
\dd s=c\,\dd\tau.
\end{equation*}

If \(\detf{\dd\pv{X}}=0\) with \(\dd\pv{X}\neq0\), the interval is null or light-like. Then
\begin{equation*}
\dd t^2=\dd\vct{r}^2
\qquad\eqtext{implies}\qquad
\dd t=\pm \|\dd\vct{r}\|,
\end{equation*}
and \((\dd\pv{X})^{-1}\) does not exist.

If \(\detf{\dd\pv{X}}<0\), the interval is space-like.

In the paravector notation, the interval classification, the proper-time formula, and the null singularity of \((\dd\pv{X})^{-1}\) are all obtained from the single algebraic identity \(\dd\pv{X}\,\conj{\dd\pv{X}}=\dd t^2-\dd\vct{r}^2\).

\subsubsection{Four-Velocity}
The four-velocity paravector is the normalized tangent to a time-like worldline:
\begin{align*}
\frac{\dd\pv{X}}{\dd s}
&=\frac{\dd t}{\dd s}+\frac{\dd\vct{r}}{\dd s}\cdot\sigv\notag\\
&=\frac{\dd t+\sigv\cdot\dd\vct{r}}{\sqrt{\dd t^2-\dd\vct{r}^2}}\notag\\
&=\frac{\dd t/\dd t+\sigv\cdot(\dd\vct{r}/\dd t)}{\sqrt{(\dd t/\dd t)^2-(\dd\vct{r}/\dd t)^2}}\notag\\
&=\frac{1+\vct{v}\cdot\sigv}{\sqrt{1-\vct{v}^2}}\notag\\
&=\gamma(1+\vct{v}\cdot\sigv).
\end{align*}
Its scalar and vector parts are therefore
\begin{equation*}
\eqtext{scalar}\!\left(\frac{\dd\pv{X}}{\dd s}\right)=\gamma,
\qquad
\eqtext{vector}\!\left(\frac{\dd\pv{X}}{\dd s}\right)=\gamma\vct{v}.
\end{equation*}
Its determinant is unity:
\begin{align*}
\detf{\frac{\dd\pv{X}}{\dd s}}
&=\gamma^2(1+\vct{v}\cdot\sigv)(1-\vct{v}\cdot\sigv)\notag\\
&=\gamma^2(1-\vct{v}^2)\notag\\
&=1.
\end{align*}
Also,
\begin{equation*}
\frac{\dd\pv{X}}{\dd t}
=\frac{\dd\pv{X}}{\dd s}\frac{\dd s}{\dd t}
=\frac{1}{\gamma}\frac{\dd\pv{X}}{\dd s}.
\end{equation*}
The four-velocity paravector is defined by
\begin{align*}
\pv{U}&=\frac{\dd\pv{X}}{\dd s}\notag\\
&=\gamma+\gamma\vct{v}\cdot\sigv\notag\\
&=U_s+\vct{U}_{\vct{v}}\cdot\sigv,\\
U_s&=\gamma,\qquad
\vct{U}_{\vct{v}}=\gamma\vct{v}.
\end{align*}
Since \(\dd s^2=\dd\pv{X}\conj{\dd\pv{X}}\) with \(\dd s>0\),
\begin{align*}
\frac{\conj{\dd\pv{X}}}{\dd s}
&=(\dd\pv{X})^{-1}\dd s\notag\\
&=\left(\frac{\dd\pv{X}}{\dd s}\right)^{-1},
\end{align*}
so \(\pv{U}^{-1}=\pv{U}^*\) and
\begin{equation*}
\detf{\pv{U}}=\pv{U}\pv{U}^*=1.
\end{equation*}
The following relation is obtained by differentiation of this normalization:
\begin{align*}
\frac{\dd}{\dd s}\detf{\pv{U}}
&=\frac{\dd\pv{U}}{\dd s}\pv{U}^*+\pv{U}\left(\frac{\dd\pv{U}}{\dd s}\right)^*\notag\\
&=\left(\frac{\dd\pv{U}}{\dd s}\pv{U}^*\right)+\left(\frac{\dd\pv{U}}{\dd s}\pv{U}^*\right)^{*\mathsf H},
\end{align*}
hence
\begin{equation*}
\eqtext{scalar}\!\left(\pv{U}^*\frac{\dd\pv{U}}{\dd s}\right)=0.
\end{equation*}
The determinant identity is recovered as
\begin{align*}
\detf{\dd\pv{X}}
&=\dd\pv{X}\conj{\dd\pv{X}}\notag\\
&=\left(\frac{\dd\pv{X}}{\dd s}\dd s\right)\left(\left(\frac{\dd\pv{X}}{\dd s}\right)^*\dd s\right)\notag\\
&=\pv{U}\pv{U}^*\,\dd s^2\notag\\
&=\detf{\pv{U}}\dd s^2\notag\\
&=\dd s^2.
\end{align*}
The proper-acceleration formulas are based on this normalization and orthogonality structure.
\noindent The following table resolves the proper-velocity paravector into its scalar and spatial components.
\begin{componenttable}[tab:components-proper-velocity]{dX/ds}
\componentrow{scalar}{\gamma}{}
\componentrow{vector}{\gamma\mathbf v}{}
\end{componenttable}


\subsubsection{Four-Acceleration}
The coordinate acceleration is defined by
\begin{equation*}
\vct{a}:=\frac{\dd\vct{v}}{\dd t}=\{a_1,a_2,a_3\}\in\R^3.
\end{equation*}
With \(\gamma=(1-\vct{v}^2)^{-1/2}\), the following relations hold:
\begin{align*}
\frac{\dd\gamma}{\dd t}&=\gamma^3(\vct{v}\cdot\vct{a}),\\
\frac{\dd\gamma}{\dd s}&=\gamma^4(\vct{v}\cdot\vct{a}).
\end{align*}
The notation \(\delv:=\partial/\partial\vct{v}\) is used for the gradient
with respect to velocity. The scalar velocity derivative needed below is
\[
\delv{}\!\left(-\frac{1}{\gamma}\right)=\gamma\vct{v}.
\]
The following derivatives are obtained by direct differentiation of the vector product \(\gamma\vct{v}\):
\begin{align*}
\frac{\dd(\gamma\vct{v})}{\dd t}
&=\gamma\vct{a}+\frac{\dd\gamma}{\dd t}\vct{v}\notag\\
&=\gamma\vct{a}+\gamma^3(\vct{v}\cdot\vct{a})\vct{v},\\
\frac{\dd(\gamma\vct{v})}{\dd s}
&=\frac{\dd(\gamma\vct{v})}{\dd t}\frac{\dd t}{\dd s}\notag\\
&=\gamma^2\vct{a}+\gamma^4(\vct{v}\cdot\vct{a})\vct{v}.
\end{align*}

The four-velocity paravector is decomposed as
\begin{align*}
\pv{U}&=U_s+\vct{U}_{\vct{v}}\cdot\sigv,\\
U_s&=\gamma,\qquad
\vct{U}_{\vct{v}}=\gamma\vct{v}.
\end{align*}
Therefore the four-acceleration is
\begin{align*}
\pv{A}
&=\frac{\dd\pv{U}}{\dd s}\notag\\
&=\frac{\dd\gamma}{\dd s}+\frac{\dd(\gamma\vct{v})}{\dd s}\cdot\sigv\notag\\
&=\gamma^4(\vct{v}\cdot\vct{a})
+\Bigl(\gamma^4(\vct{v}\cdot\vct{a})\vct{v}+\gamma^2\vct{a}\Bigr)\cdot\sigv.
\end{align*}
Equivalently,
\begin{equation*}
\pv{A}=A_s+\vct{A}_{\vct{v}}\cdot\sigv,
\qquad
A_s=\gamma^4(\vct{v}\cdot\vct{a}),
\qquad
\vct{A}_{\vct{v}}=\gamma^4(\vct{v}\cdot\vct{a})\vct{v}+\gamma^2\vct{a}.
\end{equation*}

For \(\vct{v}\neq 0\), the unit direction is defined by
\(\vct{u}=\vct{v}/\|\vct{v}\|\), with \(\vct{u}^2=1\). The following
components are then obtained from the axis functions in
Eqs.\ \eqref{eq:par-function} and \eqref{eq:prp-function}:
\begin{align*}
\eqtext{par}(\vct{a})&=(\vct{a}\cdot\vct{u})\vct{u},\\
\eqtext{prp}(\vct{a})&=\vct{a}-\eqtext{par}(\vct{a}).
\end{align*}
Then
\begin{align*}
\vct{a}^2&=\eqtext{par}(\vct{a})^2+\eqtext{prp}(\vct{a})^2,\\
(\vct{v}\cdot\vct{a})^2&=\vct{v}^2\,\eqtext{par}(\vct{a})^2,
\end{align*}
so the proper-acceleration determinant is
\begin{align*}
\detf{\pv{A}}
&=-\gamma^4\Bigl(\vct{a}^2+\gamma^2(\vct{v}\cdot\vct{a})^2\Bigr)\notag\\
&=-\gamma^4\Bigl(\eqtext{prp}(\vct{a})^2+\gamma^2\eqtext{par}(\vct{a})^2\Bigr).
\end{align*}
This is strictly space-like unless the acceleration vanishes. Its invariant magnitude \(\sqrt{-\detf{\pv{A}}}\) is the proper acceleration measured in the instantaneous rest frame.

The products with \(\pv{U}\) are
\begin{align*}
\pv{U}\pv{A}
&=(U_sA_s+\vct{U}_{\vct{v}}\cdot\vct{A}_{\vct{v}})\notag\\
&\quad+\bigl(U_s\vct{A}_{\vct{v}}+A_s\vct{U}_{\vct{v}}+i(\vct{U}_{\vct{v}}\times\vct{A}_{\vct{v}})\bigr)\cdot\sigv,\\
\pv{U}^*\pv{A}
&=(U_sA_s-\vct{U}_{\vct{v}}\cdot\vct{A}_{\vct{v}})\notag\\
&\quad+\bigl(U_s\vct{A}_{\vct{v}}-A_s\vct{U}_{\vct{v}}-i(\vct{U}_{\vct{v}}\times\vct{A}_{\vct{v}})\bigr)\cdot\sigv.
\end{align*}
Their component forms are
\begin{align*}
\pv{U}\pv{A}
&=2\gamma^5(\vct{v}\cdot\vct{a})\notag\\
&\quad+\bigl(\gamma^3\vct{a}+2\gamma^5(\vct{v}\cdot\vct{a})\vct{v}\bigr)\cdot\sigv\notag\\
&\quad+i\,\gamma^3(\vct{v}\times\vct{a})\cdot\sigv,\\
\pv{U}^*\pv{A}
&=\gamma^3\vct{a}\cdot\sigv-i\,\gamma^3(\vct{v}\times\vct{a})\cdot\sigv.
\end{align*}
In particular,
\begin{equation*}
\eqtext{scalar}(\pv{U}^*\pv{A})=0,
\end{equation*}
which is the orthogonality relation between four-velocity and four-acceleration. The scalar and vector parts of \(\pv{U}\) are recorded in Table \ref{tab:components-proper-velocity}, while the corresponding scalar and vector parts of \(\pv{U}\pv{A}\) and \(\pv{U}^{*}\pv{A}\) are compared in Tables \ref{tab:components-ua} and \ref{tab:components-ustara}.
\noindent The products \(\pv{U}\pv{A}\) and \(\pv{U}^{*}\pv{A}\) are compared in the following two tables, and the orthogonality structure of proper acceleration is thereby made explicit.
\begin{componenttable}[tab:components-ua]{\pv{U}\pv{A}}
\componentrow{real scalar}{2\gamma^{5}(\vct{v}\cdot\vct{a})}{}
\componentrow{img scalar}{0}{}
\componentrow{real vector}{\gamma^{3}\vct{a}+2\gamma^{5}(\vct{v}\cdot\vct{a})\vct{v}}{}
\componentrow{img vector}{\gamma^{3}(\vct{v}\times\vct{a})}{}
\end{componenttable}

\begin{componenttable}[tab:components-ustara]{\pv{U}^{*}\pv{A}}
\componentrow{real scalar}{0}{}
\componentrow{img scalar}{0}{}
\componentrow{real vector}{\gamma^{3}\vct{a}}{}
\componentrow{img vector}{-\gamma^{3}(\vct{v}\times\vct{a})}{}
\end{componenttable}


\subsubsection{Four-Momentum}
For a massive particle with rest mass \(m>0\), the energy--momentum paravector is produced by multiplication of four-velocity by \(m\):
\begin{equation*}
m\pv{U}=E+\vct{P}\cdot\sigv.
\end{equation*}
Its scalar part is the relativistic energy,
\begin{align*}
E
&:=\eqtext{scalar}(m\pv{U})\notag\\
&=m\gamma.
\end{align*}
If \(\vct{v}=0\), then
\begin{equation*}
E=m
\qquad \text{rest energy}
\qquad (=mc^{2}\text{ in SI units}).
\end{equation*}
The kinetic energy is therefore \(E-m\).

Its vector part is the relativistic momentum,
\begin{align*}
\vct{P}
&:=\eqtext{vector}(m\pv{U})\notag\\
&=m\gamma\vct{v}.
\end{align*}

Because \(\detf{\pv{U}}=1\), the mass shell is obtained immediately:
\begin{align*}
\detf{m\pv{U}}
&=m^{2}\notag\\
&=(E+\vct{P}\cdot\sigv)(E-\vct{P}\cdot\sigv)\notag\\
&=E^{2}-\vct{P}^{2}.
\end{align*}
Hence
\begin{equation*}
E^{2}=\vct{P}^{2}+m^{2}.
\end{equation*}
The null branch \(E^{2}=\vct{P}^{2}\), appropriate to massless motion, is also admitted by the mass-shell equation. That branch must be introduced directly as a null energy--momentum paravector: Eq.\ \(m\pv{U}=E+\vct{P}\cdot\sigv\) is applicable only to massive time-like motion because \(\dd s=0\) and \(\pv{U}=\dd\pv{X}/\dd s\) is undefined on a null worldline.

\subsubsection{Flat Paravector Coordinate Map}
The scalar and vector maps are taken to have the domains
\begin{align*}
f&:\R\times\R^3\to\R,\\
\vct{g}&:\R\times\R^3\to\R^3.
\end{align*}
A real paravector-valued map into the fixed flat space-time algebra is then
defined by
\begin{equation}
\pv{X}:=f(t,\vct{r})+\vct{g}(t,\vct{r})\cdot\sigv.
\label{eq:flat-paravector-coordinate-map}
\end{equation}
The original coordinates are mapped into flat Minkowski paravector components; invariance of the Cartesian metric coefficients under an arbitrary coordinate change is not asserted. Its differential and conjugate are
\begin{equation*}
\dd\pv{X}=\dd f+\dd\vct{g}\cdot\sigv,
\qquad
\conj{\dd\pv{X}}=\dd f-\dd\vct{g}\cdot\sigv.
\end{equation*}
Their product is
\begin{align*}
\detf{\dd\pv{X}}
&=\dd\pv{X}\,\conj{\dd\pv{X}}\notag\\
&=(\dd f+\dd\vct{g}\cdot\sigv)(\dd f-\dd\vct{g}\cdot\sigv)\notag\\
&=\dd f^2-\dd\vct{g}^2.
\end{align*}
Along a future-directed time-like path, the conditions \(\dd f>0\) and
\(\dd f^2>\dd\vct{g}^2\) are assumed. Then
\begin{align*}
\dd s
&=\sqrt{\dd f^2-\dd\vct{g}^2}\notag\\
&=\sqrt{1-\left(\frac{\dd\vct{g}}{\dd f}\right)^2}\,\dd f.
\end{align*}
The corresponding four-velocity is
\begin{align*}
\pv{U}
&=\frac{\dd\pv{X}}{\dd s}\notag\\
&=\frac{\dd f+\dd\vct{g}\cdot\sigv}{\sqrt{\dd f^2-\dd\vct{g}^2}}\notag\\
&=\frac{1+\left(\frac{\dd\vct{g}}{\dd f}\right)\cdot\sigv}{\sqrt{1-\left(\frac{\dd\vct{g}}{\dd f}\right)^2}}\notag\\
&=\frac{1+\vct{v}\cdot\sigv}{\sqrt{1-\vct{v}^2}}\notag\\
&=\gamma(1+\vct{v}\cdot\sigv),
\end{align*}
where the velocity relative to the scalar component \(f\) and its Lorentz factor are
\begin{align*}
\vct{v}
&:=\frac{\dd\vct{g}}{\dd f}\notag\\
&=\frac{\dd\vct{g}/\dd t}{\dd f/\dd t}\notag\\
&=\frac{\dd\vct{g}/\dd s}{\dd f/\dd s},\\
\gamma
&:=\frac{\dd f}{\dd s}\notag\\
&=\frac{\dd f}{\sqrt{\dd f^2-\dd\vct{g}^2}}\notag\\
&=\frac{\dd f/\dd t}{\dd s/\dd t}\notag\\
&=\frac{1}{\sqrt{1-\vct{v}^2}}.
\end{align*}
Thus \(\vct{v}^2<1\). For a massive particle, the energy--momentum relation remains unchanged:
\begin{equation*}
m\pv{U}=E+\vct{P}\cdot\sigv.
\end{equation*}
The coordinate acceleration relative to \(f\) is defined, in direct analogy
with \(\vct{a}=\dd\vct{v}/\dd t\) above, by
\begin{align*}
\vct{a}
&:=\frac{\dd\vct{v}}{\dd f}\notag\\
&=\frac{\dd^2\vct{g}}{\dd f^2}.
\end{align*}
Because \(\dd/\dd s=\gamma\,\dd/\dd f\), the four-acceleration has the same form as before:
\begin{equation*}
\frac{\dd\pv{U}}{\dd s}
=\gamma^4(\vct{v}\cdot\vct{a})
+\bigl(\gamma^4(\vct{v}\cdot\vct{a})\vct{v}+\gamma^2\vct{a}\bigr)\cdot\sigv.
\end{equation*}
The commutator of four-velocity and four-acceleration is
\begin{align*}
\com{\pv{U}}{\dfrac{\dd\pv{U}}{\dd s}}
&=\com{\pv{U}}{\dfrac{\dd\pv{U}}{\dd f}}\frac{\dd f}{\dd s}\notag\\
&=\gamma\com{\pv{U}}{\dfrac{\dd\pv{U}}{\dd f}}\notag\\
&=2\gamma^3(\vct{v}\times\vct{a})\cdot(i\sigv).
\end{align*}
For a massive free particle in these flat target components,
\begin{equation*}
m\frac{\dd\pv{U}}{\dd s}=0,
\end{equation*}
which is equivalently the straight-worldline equation
\begin{equation*}
\frac{\dd^2\pv{X}}{\dd s^2}=0.
\end{equation*}
After division by the nonzero factor \(\gamma\), its scalar and vector parts are
\begin{align*}
\gamma^3(\vct{v}\cdot\vct{a})&=0,\\
\gamma^3(\vct{v}\cdot\vct{a})\vct{v}+\gamma\vct{a}&=0.
\end{align*}
For \(\vct{v}^2<1\), one has \(\gamma>0\). The relation
\(\vct{v}\cdot\vct{a}=0\) is therefore obtained from the scalar part, and
\(\vct{a}=0\) is then obtained from the vector part. Hence
\begin{equation*}
\frac{\dd^2\vct{g}}{\dd f^2}=0,
\end{equation*}
and along an inertial worldline
\begin{equation*}
\vct{g}=\vct{c}f+\vct{c}',
\qquad
\vct{v}=\frac{\dd\vct{g}}{\dd f}=\vct{c},
\end{equation*}
so inertial motion is linear in the flat target components \((f,\vct{g})\) with constant velocity. If the same straight worldline is described in arbitrary curvilinear coordinates, the corresponding connection terms are generally contained in its coordinate equation; no curved-space geodesic equation is asserted.

\subsubsection{Lagrangian Mechanics and Conserved Quantities}
For a particle trajectory written in terms of time, position, and velocity,
the action is taken to be
\begin{equation*}
S=\int \mathcal L(t,\vct{r},\vct{v})\,\dd t.
\end{equation*}
The Euler--Lagrange equation is then
\begin{align*}
\left(\gradv-\frac{\dd}{\dd t}\delv{}\right)\mathcal L&=0,\\
\frac{\dd}{\dd t}\delv{}\mathcal L&=\gradv\mathcal L.
\end{align*}
The standard conserved quantities are then obtained immediately. The canonical momentum is
\begin{equation*}
\delv{}\mathcal L,
\end{equation*}
so if \(\gradv\mathcal L=0\), then
\begin{equation*}
\frac{\dd}{\dd t}\delv{}\mathcal L=0,
\end{equation*}
and the canonical momentum is constant.

The Hamiltonian is defined by
\begin{align*}
\mathcal H
&=\vct{v}\cdot(\delv{}\mathcal L)-\mathcal L\notag\\
&=(\vct{v}\cdot\delv{}-1)\mathcal L.
\end{align*}
Its time derivative is
\begin{align*}
\frac{\dd\mathcal H}{\dd t}
&=\frac{\dd}{\dd t}(\delv{}\mathcal L)\cdot\vct{v}
+(\delv{}\mathcal L)\cdot\frac{\dd\vct{v}}{\dd t}
-\frac{\dd\mathcal L}{\dd t}\notag\\
&=(\gradv\mathcal L)\cdot\vct{v}
+(\delv{}\mathcal L)\cdot\frac{\dd\vct{v}}{\dd t}
-\frac{\dd\mathcal L}{\dd t}\notag\\
&=(\gradv\mathcal L)\cdot\vct{v}
+(\delv{}\mathcal L)\cdot\frac{\dd\vct{v}}{\dd t}\notag\\
&\qquad-\left(
\partial_t\mathcal L
+(\gradv\mathcal L)\cdot\vct{v}
+(\delv{}\mathcal L)\cdot\frac{\dd\vct{v}}{\dd t}
\right)\notag\\
&=-\partial_t\mathcal L.
\end{align*}
Therefore, if \(\partial_t\mathcal L=0\), then
\begin{equation*}
\frac{\dd\mathcal H}{\dd t}=0,
\end{equation*}
so the Hamiltonian is constant. These identities are applied directly to the free-particle and minimally coupled electromagnetic examples.

\subsubsection{Action of a Free Particle}
The free-particle action is the mass-weighted proper time along the world line. With \(\tau=\int \dd s\), \(\pv{U}=\dd\pv{X}/\dd s\), and \(\gamma=\dd t/\dd s\),
\begin{align*}
S
=\int \mathcal L\,\dd t
&=-m\tau\notag\\
&=-m\int \sqrt{\detf{\dd\pv{X}}}\notag\\
&=-m\int \dd s\notag\\
&=-m\int \sqrt{\detf{\frac{\dd\pv{X}}{\dd t}}}\,\dd t\notag\\
&=-m\int \frac{\dd s}{\dd t}\,\dd t\notag\\
&=-m\int \sqrt{\detf{\frac{\pv{U}}{\gamma}}}\,\dd t\notag\\
&=-m\int \frac{1}{\gamma}\,\dd t\notag\\
&=-m\int \sqrt{1-\vct{v}^2}\,\dd t.
\end{align*}
Therefore
\begin{equation*}
\mathcal L=-\frac{m}{\gamma}.
\end{equation*}
The canonical momentum is obtained by differentiation with respect to velocity:
\begin{align*}
\delv{} \mathcal L
&=m\gamma\vct{v}\notag\\
&=\vct{P},\\
(\delv{}\cdot\sigv)\mathcal L
&=\vct{P}\cdot\sigv.
\end{align*}
The Hamiltonian is
\begin{align*}
\mathcal H
&=(\delv{}\mathcal L)\cdot\vct{v}-\mathcal L\notag\\
&=m\gamma\vct{v}^2+\frac{m}{\gamma}\notag\\
&=m\gamma\left(\vct{v}^2+\frac{1}{\gamma^2}\right)\notag\\
&=m\gamma\notag\\
&=E.
\end{align*}
Because the free Lagrangian has no explicit time or position dependence,
\begin{equation*}
\partial_t\mathcal L=0,
\qquad
\gradv\mathcal L=0,
\end{equation*}
so the conserved quantities are the relativistic energy \(E=m\gamma\) and momentum \(\vct{P}=m\gamma\vct{v}\).

\subsubsection{Non-Relativistic Approximation}
For the non-relativistic approximation, the free Lagrangian is expanded
directly:
\begin{align*}
\mathcal L
&=-\frac{m}{\gamma}\notag\\
&=-m+\frac12 m\vct{v}^2+O(\vct{v}^4).
\end{align*}
The canonical momentum is then reduced to
\begin{align*}
\delv{}\mathcal L
&=\frac{\partial}{\partial \vct{v}}\left(\frac12 m\vct{v}^2+O(\vct{v}^4)\right)\notag\\
&=m\vct{v}+O(\vct{v}^3)\notag\\
&\approx m\vct{v}.
\end{align*}
The total energy is the kinetic term plus the rest-mass term:
\begin{align*}
\mathcal H
&=\bigl(m\vct{v}+O(\vct{v}^3)\bigr)\cdot\vct{v}
+m-\frac12m\vct{v}^2-O(\vct{v}^4)\notag\\
&=m+\frac12m\vct{v}^2+O(\vct{v}^4)\notag\\
&\approx m+\frac12m\vct{v}^2.
\end{align*}
Equivalently, the low-velocity series is
\begin{align*}
\mathcal L&=-m+\frac12 m\vct{v}^2+O(\vct{v}^4),\notag \\
\vct{P}&=m\vct{v}+O(\vct{v}^{3}),\notag \\
\mathcal H&=m+\frac12 m\vct{v}^{2}+O(\vct{v}^{4}).
\end{align*}
The Newtonian kinetic term together with the additive rest-energy constant is thereby shown to be contained in the relativistic action. Classical mechanics is therefore obtained as the approximation, not used as the starting point.

\subsubsection{Flat Pullback Metric and Diagonalization}
A flat metric in the original coordinate tuple is induced by the map \(\pv{X}=f+\vct{g}\cdot\sigv\):
\[
\vct{x}=(t,r_1,r_2,r_3)^T,
\qquad
\dd\vct{x}=(\dd t,\dd r_1,\dd r_2,\dd r_3)^T.
\]
The transpose operation was defined in Eq.\ \eqref{eq:matrix-transpose}.
The symmetric pullback matrix \(\Qmat(\vct{x})\) is the unique matrix defined
by expansion of \(\dd f\) and \(\dd\vct{g}\) in \(\dd\vct{x}\):
\begin{equation}
\dd f^2-\dd\vct{g}^2
=:\dd\vct{x}^T\Qmat(\vct{x})\,\dd\vct{x}.
\label{eq:pullback-matrix-function}
\end{equation}
Consequently,
\begin{equation*}
\dd s^2=\dd\vct{x}^T\Qmat(\vct{x})\,\dd\vct{x}.
\end{equation*}
If the map is locally invertible, the signature \((+,-,-,-)\) is preserved by Sylvester's law. Position-dependent entries of \(\Qmat\) can be produced by curvilinear coordinates even though the target geometry remains flat; curvature is not represented by these entries alone. If the map fails to be locally invertible, the pullback may instead be degenerate.

At every point where the map is locally invertible, the ordered eigenvalue
function is defined by
\begin{equation}
\lambda(n):=\eqtext{ordered eigenvalue }n\eqtext{ of }\Qmat,
\qquad n\in\{1,2,3,4\}.
\label{eq:metric-eigenvalue-function}
\end{equation}
An orthogonal decomposition is then obtained for the real symmetric matrix:
\[
\Qmat=\Nmat\eqtext{diag}(\lambda(1),\lambda(2),\lambda(3),\lambda(4))\Nmat^T,
\qquad
\Nmat^T\Nmat=\I.
\]
The eigenvalues are ordered so that
\[
\lambda(1)>0,
\qquad
\lambda(2)<0,
\qquad
\lambda(3)<0,
\qquad
\lambda(4)<0.
\]
The corresponding direction function is defined by
\begin{equation}
\vct{u}(n):=\eqtext{column }n\eqtext{ of }\Nmat.
\label{eq:metric-direction-function}
\end{equation}
The local differential directions are defined by
\begin{equation}
\omega(n):=\sqrt{|\lambda(n)|}\,\vct{u}(n)^T\dd\vct{x}.
\label{eq:local-differential-directions}
\end{equation}
The interval is consequently
\begin{equation*}
\dd s^2
=\omega(1)^2-\omega(2)^2-\omega(3)^2-\omega(4)^2.
\end{equation*}
The three minus signs cannot be dropped: a Lorentzian interval is not
turned into a Euclidean norm by diagonalization. When \(\Qmat\) varies with
position, the local differential directions \(\omega(n)\) need not be
derivatives of globally defined coordinates.
